%% notes-tex/multimodal/043-training-plan/sections/5-grid.tex
%% SOURCE: measured step times from `sacct -j 6136,6137,6150,6152 -X --format=Elapsed` (GilaHyper,
%% 2026-10-23): 6136 03:06:30, 6137 07:12:38, 6150 03:03:59, 6152 02:33:59, each over 317,300
%% optimizer steps at B = 8 (design review tables/macros.tex); everything else in this file is
%% arithmetic on those four numbers, the step counts of Secs. 2 and 3, and the machine factors of
%% Table tab:cost, which are planning values, not measurements. Machine and queue facts from
%% memory alcf-polaris-login (qstat -Qf 2026-10-17; allocations 2026-10-16), the owner's machine
%% list of 2026-10-23, and memory delta-never-preempt-partition.

\section{The grid\secstatus{todo}}
\label{sec:grid}

\subsection{The cost model\secstatus{todo}}
\label{sec:cost}

Every hour in this section is an estimate from four measured runs and a machine factor
(Table~\ref{tab:cost}). A run's estimate is its optimizer steps times the measured seconds per
step of its closest overnight run, times the factor of the machine. The seconds per step are
whole-job elapsed time over steps, so they include evaluation and startup. The factors are
planning values: A100 at 1.0, A40 at 1.5 and an Aurora tile at 2.0 times the RTX 6000 Ada time.
None is measured for this trainer; the first smoke on each machine replaces its factor, and
every \texttt{-{}-time} and PBS walltime below is recomputed from that measurement before
submission. Requested time is about 1.5 times the estimate (the prepare-experiment gate's
margin). Route B and the 4{,}096-mark cap are budgeted at twice the reference's time, a guess
that smoke G6 replaces.

%% SOURCE: sacct elapsed times of jobs 6136, 6137, 6150, 6152 (header of this file).
\begin{table}[htbp]
\centering
\footnotesize
\setlength{\tabcolsep}{3pt}
\caption[]{\textbf{Measured seconds per optimizer step and the planning factors.} Top:
GilaHyper, one RTX 6000 Ada card, batch 8, elapsed time over 317{,}300 steps. Bottom: factors
on that time, unmeasured.}
\label{tab:cost}
\begin{tabular}{>{\raggedright\arraybackslash}p{62mm}>{\raggedright\arraybackslash}p{20mm}>{\raggedleft\arraybackslash}p{24mm}>{\raggedleft\arraybackslash}p{24mm}>{\raggedright\arraybackslash}p{36mm}}
\toprule
Run & Slurm job & Elapsed, h:mm & Seconds per step, s & Used for \\
\midrule
\texttt{joint\_9}, records sampler, HL-Gauss & 6136 & 3:07 & 0.0353 & records-sampler arms \\
\texttt{joint\_9}, records sampler, pinball & 6150 & 3:04 & 0.0348 & (agrees with 6136) \\
\texttt{joint\_9}, balanced sampler, HL-Gauss & 6137 & 7:13 & 0.0818 & every other stage-1 arm; segregant fine-tunes \\
Kemmeren alone & 6152 & 2:34 & 0.0291 & Kemmeren fine-tunes and controls \\
\midrule
Machine factor, A100 40\,GB (Polaris; Delta A100) & -- & -- & $\times 1.0$ & planning value \\
Machine factor, A40 48\,GB (Delta) & -- & -- & $\times 1.5$ & planning value \\
Machine factor, Max 1550 tile (Aurora) & -- & -- & $\times 2.0$ & planning value \\
\bottomrule
\end{tabular}
\end{table}

With $k_1 \approx 353{,}150$, a stage-1 reference run is $353{,}150 \times 0.0818$\,s $= 8.0$\,h
on an Ada card, a records-sampler run 3.5\,h; a Kemmeren fine-tune of 31{,}710 steps is 0.26\,h;
a Kemmeren $\mathrm{alone}(k_{\mathrm{tot}})$ of 384{,}860 steps is 3.1\,h; a segregant
fine-tune 0.72\,h and a segregant $\mathrm{alone}(k_{\mathrm{tot}})$ 8.75\,h.

\subsection{(a) GilaHyper smokes, today\secstatus{todo}}
\label{sec:gilahyper}

Two cards, one smoke per card at a time, chained \texttt{afterany}. Each smoke runs on the real
prepared store with the real logger and must pass the gate of Sec.~\ref{sec:launch} before any
run it enables is submitted. Every smoke but G7 waits on code that is not in the trainer at
commit \texttt{f1cf262d2}; the last column names it.

%% SOURCE: arithmetic on Table tab:cost (Ada, factor 1.0); 7,063 steps per epoch of `all` (estimate, Sec. 2.1).
\begin{table}[htbp]
\centering
\footnotesize
\setlength{\tabcolsep}{3pt}
\caption[]{\textbf{Block (a): GilaHyper smokes.} One RTX 6000 Ada card per job, seed 0,
batch 8. Estimated hours from Table~\ref{tab:cost}; requested is the slurm
\texttt{-{}-time}.}
\label{tab:gh}
\begin{tabular}{>{\raggedright\arraybackslash}p{8mm}>{\raggedright\arraybackslash}p{44mm}>{\raggedright\arraybackslash}p{20mm}>{\raggedleft\arraybackslash}p{14mm}>{\raggedleft\arraybackslash}p{14mm}>{\raggedright\arraybackslash}p{30mm}>{\raggedright\arraybackslash}p{36mm}}
\toprule
Id & Run & Length & Estimated, card-h & Requested, h:mm & Enables & Needs first \\
\midrule
G1 & \texttt{all}, reference arm (balanced, HL-Gauss, equal, plain residual, mask) & 2 epochs & 0.4 & 1:00 & P0, the Delta request, the Polaris stage-1 job & specs for Hu, Jackson, Nadal-Ribelles; \texttt{conf/arm/all.yaml}; prepare of \texttt{all}; residual switch \\
G2 & fine-tune path: 6136's window-end checkpoint (\texttt{joint\_9}) onto Kemmeren & 2 epochs (302 steps) & 0.1 & 0:30 & D5, D6, Polaris stage 2 & checkpoint load mapping dataset rows by name; $\mathrm{lr}_2$, $\mathcal{F}$; standardization sha256 check \\
G3 & readout gather (Eq.~\ref{eq:gather}) & 1 epoch & 0.2 & 0:45 & C5 & graph-placement switch \\
G4 & no graph & 1 epoch & 0.2 & 0:45 & C6 & graph-placement switch \\
G5 & background marks (Eq.~\ref{eq:bgmarks}) & 1 epoch & 0.2 & 0:45 & C7 & background registry, marks, $\mathbf{E}^{\mathrm{orig}}$ \\
G6 & memory probe: route B, and uncapped marks, on the largest-mark batches & 50 steps each & 0.1 & 0:30 & C8, C9, and the factor of 2 & route B in the trainer \\
G7 & Kemmeren $\mathrm{alone}(k)$, seeds 0 to 5 & 31{,}710 steps each & 1.6 & 3:00 (6 jobs of 0:30) & the $\Delta^{\mathrm{ft}}_{k}$ controls & nothing \\
\midrule
\multicolumn{3}{l}{Block total} & 2.8 & 7:15 & & \\
\bottomrule
\end{tabular}
\end{table}

G2 loads a \texttt{joint\_9} checkpoint into an \texttt{all} model on purpose: the two have
different dataset vocabularies, so the smoke proves that rows of $\mathbf{E}^{\mathrm{data}}$
and $\mathbf{b}_s$ are matched by dataset name, not by position.

\subsection{(b) Delta, eight individual jobs (owner go required)\secstatus{todo}}
\label{sec:delta}

Eight single-card jobs on Delta's A40 partition, never the preempt partition, submitted only
after the owner's go. Their role is an early read of the whole two-stage chain at two seeds
while the Polaris job queues, and a cross-machine check: seeds 1 and 2 repeat inside the
Polaris grid. D7 runs first and doubles as the A40 timing probe; D1 to D4 are submitted after
its first evaluation, with \texttt{-{}-time} recomputed from its measured seconds per step.

%% SOURCE: arithmetic on Table tab:cost with the A40 factor 1.5.
\begin{table}[htbp]
\centering
\footnotesize
\setlength{\tabcolsep}{3pt}
\caption[]{\textbf{Block (b): Delta, eight jobs on single A40 cards.} Stage-1 runs use the
reference configuration except where the arm says otherwise. D5 and D6 depend on D1 and D2
with \texttt{afterany} and raise if the checkpoint is missing.}
\label{tab:delta}
\begin{tabular}{>{\raggedright\arraybackslash}p{8mm}>{\raggedright\arraybackslash}p{50mm}>{\raggedleft\arraybackslash}p{10mm}>{\raggedright\arraybackslash}p{26mm}>{\raggedleft\arraybackslash}p{10mm}>{\raggedleft\arraybackslash}p{16mm}>{\raggedleft\arraybackslash}p{16mm}>{\raggedright\arraybackslash}p{18mm}}
\toprule
Id & Arm & Seed & Length & Cards, n & Estimated, card-h & Requested, h:mm & Dependency \\
\midrule
D1 & stage 1, reference (plain residual) & 1 & 50 epochs & 1 & 12.0 & 18:00 & after D7 read \\
D2 & stage 1, reference (plain residual) & 2 & 50 epochs & 1 & 12.0 & 18:00 & after D7 read \\
D3 & stage 1, ReZero control (C1) & 1 & 50 epochs & 1 & 12.0 & 18:00 & after D7 read \\
D4 & stage 1, ReZero control (C1) & 2 & 50 epochs & 1 & 12.0 & 18:00 & after D7 read \\
D5 & stage 2, $\mathrm{ft}(k)$ on Kemmeren from D1 & 1 & 210 epochs & 1 & 0.4 & 1:00 & afterany D1 \\
D6 & stage 2, $\mathrm{ft}(k)$ on Kemmeren from D2 & 2 & 210 epochs & 1 & 0.4 & 1:00 & afterany D2 \\
D7 & Kemmeren $\mathrm{alone}(k_{\mathrm{tot}})$, 384{,}860 steps & 1 & 2{,}549 epochs & 1 & 4.7 & 7:00 & none \\
D8 & Kemmeren $\mathrm{alone}(k_{\mathrm{tot}})$, 384{,}860 steps & 2 & 2{,}549 epochs & 1 & 4.7 & 7:00 & none \\
\midrule
\multicolumn{4}{l}{Block total} & 8 & 58.2 & 88:00 & \\
\bottomrule
\end{tabular}
\end{table}

The $\mathrm{alone}(k)$ controls of seeds 1 and 2 come from G7. Two seeds decide nothing under
the rule; the block's outputs are the A40 factor, a working chain, and the first paired
values of $\Delta^{\mathrm{ft}}_{\mathrm{tot}}$ and $\Delta^{\mathrm{ft}}_{k}$.

\subsection{(c) Polaris\secstatus{todo}}
\label{sec:polaris}

Polaris holds the six-seed stage-1 grid. Each node has four A100 40\,GB cards and each run
takes one card, so a packed job runs four runs per node: one process per card
(\texttt{CUDA\_VISIBLE\_DEVICES}), drawing from a task list ordered longest first, each task its
own W\&B run, a failed task leaving the others running. The queue is \texttt{capacity} (1 to 4
nodes, 168\,h, one job per project at a time). The \texttt{preemptable} queue is not used: the
trainer has no resume, so a preemption loses the run. The 1-hour \texttt{debug} queue takes the
timing probe and the short fine-tunes.

%% SOURCE: arithmetic on Table tab:cost (A100 factor 1.0) and the arms of Secs. 2.2 to 2.7.
\begingroup
\footnotesize
\setlength{\tabcolsep}{3pt}
\begin{longtable}{>{\raggedright\arraybackslash}p{9mm}>{\raggedright\arraybackslash}p{70mm}>{\raggedleft\arraybackslash}p{12mm}>{\raggedleft\arraybackslash}p{22mm}>{\raggedleft\arraybackslash}p{20mm}>{\raggedright\arraybackslash}p{24mm}}
\caption[]{\textbf{Block (c), the Polaris stage-1 job: arms, seeds 0 to 5.} One change from
the reference per contrast. Estimated hours per run on one A100 card; card-hours are runs times
hours.}\label{tab:polaris}\\
\toprule
Id & Arm (change from the reference) & Runs, n & Estimated per run, card-h & Estimated, card-h & Control \\
\midrule
\endfirsthead
\toprule
Id & Arm (change from the reference) & Runs, n & Estimated per run, card-h & Estimated, card-h & Control \\
\midrule
\endhead
R & reference: \texttt{all}, balanced, HL-Gauss, equal, plain residual, mask at layer 2, no background marks, route A, cap 1{,}024 & 6 & 8.0 & 48.0 & -- \\
C1 & ReZero residual (\deq{9}) & 6 & 8.0 & 48.0 & R \\
C2 & records sampler, equal weighting & 6 & 3.5 & 21.0 & R \\
C3 & records sampler, record weighting & 6 & 3.5 & 21.0 & C2 \\
C4 & sampler $\alpha = 0.5$ (Eq.~\ref{eq:sampler}), equal weighting & 6 & 8.0 & 48.0 & R \\
C5 & readout gather, no mask (Eq.~\ref{eq:gather}) & 6 & 8.0 & 48.0 & R \\
C6 & no graph & 6 & 8.0 & 48.0 & R \\
C7 & background marks (Eq.~\ref{eq:bgmarks}) & 6 & 8.0 & 48.0 & R \\
C8 & route B for the isolates, uncapped & 6 & 16.0 & 96.0 & R \\
C9 & mark cap 4{,}096 & 6 & 16.0 & 96.0 & R \\
C10 & pinball loss & 6 & 8.0 & 48.0 & R \\
C11 & \texttt{joint\_9} at $k_1$ steps (without Hu, Jackson, Nadal-Ribelles) & 6 & 8.0 & 48.0 & R \\
K & Kemmeren $\mathrm{alone}(k_{\mathrm{tot}})$, the stage-2 control & 6 & 3.1 & 18.6 & -- \\
\midrule
 & Job total & 78 & & 636.6 & \\
\bottomrule
\end{longtable}
\endgroup

On 16 cards the job's 636.6 card-hours take about 39.8\,h, which is 159.2 node-hours; the
longest single runs (16\,h) fit well inside that. Every stage-1 run saves its window-end
checkpoint, so stage 2 can start from whichever arm the rule names without rerunning stage 1.

%% SOURCE: arithmetic on Tables tab:cost and tab:polaris; allocation 3,884 node-hours
%% (GenesisHackathonOct26, sbank 2026-10-16).
\begin{table}[htbp]
\centering
\footnotesize
\setlength{\tabcolsep}{3pt}
\caption[]{\textbf{Block (c): Polaris jobs and node-hours.} Estimated node-hours are nodes
times the estimated wall time; requested is nodes times the PBS walltime. The allocation is
3{,}884 node-hours on \texttt{GenesisHackathonOct26}.}
\label{tab:polarisjobs}
\begin{tabular}{>{\raggedright\arraybackslash}p{9mm}>{\raggedright\arraybackslash}p{58mm}>{\raggedright\arraybackslash}p{18mm}>{\raggedleft\arraybackslash}p{10mm}>{\raggedleft\arraybackslash}p{17mm}>{\raggedleft\arraybackslash}p{17mm}>{\raggedleft\arraybackslash}p{17mm}}
\toprule
Id & Job & Queue & Nodes, n & Walltime requested, h & Estimated, node-h & Requested, node-h \\
\midrule
P0 & timing probe: R for 2 epochs on four cards (seeds 0 to 3), C8 memory, the fine-tune path & debug & 1 & 1 & 0.5 & 1 \\
P1 & stage-1 grid, Table~\ref{tab:polaris} & capacity & 4 & 60 & 159.2 & 240 \\
P2 & stage 2 from R, seeds 0 to 5: Kemmeren $\mathrm{ft}(k)$; singles $\mathrm{ft}(1{,}100)$ and $\mathrm{alone}(1{,}100)$ & debug & 1 & 1 & 0.6 & 1 \\
\midrule
 & Block total & & & & 160.3 & 242 \\
\bottomrule
\end{tabular}
\end{table}

\subsection{(d) Aurora\secstatus{todo}}
\label{sec:aurora}

Aurora runs only after the trainer runs on Intel XPU and its results agree with Polaris'. Its
torch 2.11 XPU environment is live and PyG works on XPU (17 of 17 core checks, job 8938797);
what the 043 trainer still needs:

\begin{enumerate}[leftmargin=*,itemsep=2pt]
\item \textbf{Device.} \texttt{TrainConfig043.device} admits \texttt{xpu}; device selection,
\texttt{pin\_memory} and the bf16 autocast (\texttt{device\_type="xpu"}) follow it. At commit
\texttt{f1cf262d2} the trainer knows only \texttt{cuda} and \texttt{cpu}.
\item \textbf{The \texttt{torch\_scatter} import chain.} \file{cell_model_043.py} imports
\file{equivariant_cell_graph_transformer.py}, whose package import reaches
\file{torchcell/models/__init__.py} and \file{deep_set.py}, which imports \texttt{torch\_scatter};
there is no XPU build of it. The candidate fix, \texttt{torch\_geometric.utils.scatter} or a
lazy import, is untested.
\item \textbf{Attention kernels.} The masked trunk layer and the operator call
\texttt{scaled\_dot\_product\_attention} with masks; its output on XPU is compared with CUDA on
a fixed batch and fixed weights at bf16 tolerance, and its speed is measured.
\item \textbf{The launcher.} A PBS script with \texttt{ZE\_FLAT\_DEVICE\_HIERARCHY=FLAT}, one
process per tile pinned by \texttt{ZE\_AFFINITY\_MASK}, twelve per node, and W\&B online through
the ALCF proxy.
\item \textbf{Data.} The prepared store, the frozen vector tables, the split file and the
standardization records on \texttt{/flare}, each verified by sha256 against the GilaHyper copy
(the transfer starts 2026-10-23).
\item \textbf{Parity.} Job A1 runs R and C1, seeds 0 to 5, on twelve tiles. Aurora results enter
a decision only if, per dataset of Table~\ref{tab:ladder}, the paired Aurora-minus-Polaris
window-mean difference has a 95\% interval covering zero and a mean below 0.01 in magnitude, half
the rule's margin (a proposed threshold).
\end{enumerate}

%% SOURCE: arithmetic on Table tab:cost with the tile factor 2.0; allocation 18,468 node-hours
%% (sbank 2026-10-16); queue limits from qstat -Qf 2026-10-17.
\begin{table}[htbp]
\centering
\footnotesize
\setlength{\tabcolsep}{3pt}
\caption[]{\textbf{Block (d): Aurora jobs and node-hours.} Twelve tiles per node, one run per
tile. Estimated tile-hours per run: stage-1 reference 16.0; Kemmeren fine-tune 0.5; segregant
fine-tune and $\mathrm{alone}(k)$ 1.4; segregant $\mathrm{alone}(k_{\mathrm{tot}})$ 17.5. The
allocation is 18{,}468 node-hours.}
\label{tab:aurora}
\begin{tabular}{>{\raggedright\arraybackslash}p{9mm}>{\raggedright\arraybackslash}p{62mm}>{\raggedright\arraybackslash}p{18mm}>{\raggedleft\arraybackslash}p{10mm}>{\raggedleft\arraybackslash}p{15mm}>{\raggedleft\arraybackslash}p{15mm}>{\raggedleft\arraybackslash}p{15mm}}
\toprule
Id & Job & Queue & Nodes, n & Walltime requested, h & Estimated, node-h & Requested, node-h \\
\midrule
A0 & XPU port tests: CUDA parity of a fixed batch, one-tile smoke, twelve-tile launcher (4 jobs) & debug & 1 & 1 each & 2 & 4 \\
A1 & parity: R and C1, seeds 0 to 5, 50 epochs & capacity & 1 & 24 & 16 & 24 \\
A2 & stage-2 grid from Polaris R checkpoints: four targets $\times$ $\mathrm{lr}_2 \in \{3, 1, 0.3\} \times 10^{-4}$ $\times$ three policies $\mathcal{F}$ $\times$ 6 seeds (216 runs), plus the segregant and singles controls (30 runs); 406 tile-h & capacity & 2 & 27 & 35 & 54 \\
A3 & longer stage 1: R and C1 at 200 epochs, seeds 0 to 5 (owner decision 19) & capacity & 1 & 96 & 64 & 96 \\
\midrule
 & Block total & & & & 117 & 178 \\
\bottomrule
\end{tabular}
\end{table}

A2 compares its Kemmeren fine-tunes with the Polaris Kemmeren controls (K and G7), a
cross-machine comparison admitted only after A1 passes. On two nodes the twelve segregant
$\mathrm{alone}(k_{\mathrm{tot}})$ runs (17.5\,h each) occupy twelve tiles, and the other twelve
tiles finish the remaining 196 tile-hours in about 16.3\,h, so the wall is about 17.5\,h.

\subsection{Totals and the order of unblocking\secstatus{todo}}
\label{sec:totals}

%% SOURCE: block totals of Tables tab:gh, tab:delta, tab:polarisjobs, tab:aurora.
\begin{table}[htbp]
\centering
\footnotesize
\setlength{\tabcolsep}{3pt}
\caption[]{\textbf{Budget per machine.} GilaHyper and Delta in card-hours, Polaris and
Aurora in node-hours, the unit each allocation is charged in. Share is requested over the
allocation.}
\label{tab:totals}
\begin{tabular}{>{\raggedright\arraybackslash}p{20mm}>{\raggedright\arraybackslash}p{30mm}>{\raggedleft\arraybackslash}p{18mm}>{\raggedleft\arraybackslash}p{18mm}>{\raggedleft\arraybackslash}p{24mm}>{\raggedleft\arraybackslash}p{20mm}>{\raggedright\arraybackslash}p{24mm}}
\toprule
Machine & Unit & Estimated & Requested & Allocation & Share of allocation, \% & Gate \\
\midrule
GilaHyper & card-h & 2.8 & 7.3 & two cards, owner's grant & -- & none \\
Delta & card-h & 58.2 & 88.0 & owner's account & -- & owner's go \\
Polaris & node-h & 160.3 & 242 & 3{,}884 & 6.2 & G1, P0 \\
Aurora & node-h & 117 & 178 & 18{,}468 & 1.0 & A0, A1 \\
\bottomrule
\end{tabular}
\end{table}

\begin{enumerate}[leftmargin=*,itemsep=2pt]
\item \textbf{Code}, by the trainer's owner: the three dataset specs, the residual, graph and
background switches, the fine-tune entry point, route B, \file{conf/decision_rule.yaml}, the
Ohya probe script, the XPU items of Sec.~\ref{sec:aurora}.
\item \textbf{GilaHyper G1 to G7} (Table~\ref{tab:gh}). G1 passing the gate unblocks P0 and the
Delta request; G2 unblocks every fine-tune; G6 fixes the budget of C8 and C9.
\item \textbf{Polaris P0} replaces the A100 factor and sets P1's walltime; \textbf{P1} then runs
the stage-1 grid and the Kemmeren $\mathrm{alone}(k_{\mathrm{tot}})$ controls.
\item \textbf{Delta}, after the owner's go and in parallel with P1: D7 and D8 first, D1 to D4 on
D7's measured factor, D5 and D6 behind D1 and D2.
\item \textbf{Aurora A0 then A1}, in parallel with P1; A1 compares with P1's R and C1.
\item \textbf{P1's result under the rule} names the stage-1 arm stage 2 starts from (R unless a
contrast wins). \textbf{P2} runs from its checkpoints; \textbf{A2}, if A1 passed, runs the
fine-tuning grid; \textbf{A3} waits for the owner.
\item \textbf{End of campaign}: fold 0 and the three doubles statements of
Sec.~\ref{sec:ftdoubles} are read once.
\end{enumerate}
